% translation-source: docs/user/latex/chapters/10-effets.tex; sha256: 631ea3a63b0b3d80150e1fbea4a21e8e3cbb4900e4f11c5a37a7ef55d7ce884c
\chapter{Effects}
\section{Choose an effect}
\begin{longtable}{@{}L{3.6cm}L{2.8cm}L{7.2cm}@{}}
\toprule Control & Range & Intention\\\midrule\endhead
Texture & $-100$ to $+100$ & Soften or enhance small details.\\
Clarity & $-100$ to $+100$ & Change structural presence and local contrast.\\
Dehaze & $-100$ to $+100$ & Change apparent haze and tonal separation.\\
Grain & 0 to 100 & Add grain; this is not noise reduction.\\
Gaussian blur & 0 to 20 & Soften globally; not a complete depth-of-field simulation.\\
\bottomrule
\end{longtable}
All these controls have zero as reference. Negative texture or clarity is
not equivalent to the dedicated noise-reduction module.

\section{Check artifacts}
Inspect edges, sky transitions and fine textures after a strong correction.
Increased clarity can make noise more visible; excessive smoothing can make
surfaces look artificial. Also check delivery size: detail that looks pleasant
at high magnification may be excessive in fit-to-view.

\paragraph{Example.}
On a hazy landscape, start with a small dehaze correction, then rebalance
light if needed. Compare using \touche{Ctrl + Shift + Y} and inspect sky/building
boundaries at \touche{Ctrl + 1}.
\gestesreglages
